%Data institusi

%\input{Dozentdaten}

%\renewcommand{\fachbereich}{Fachbereich}

%\renewcommand{\dozent}{Prof. Dr. . }

%Data ujian

\renewcommand{\klausurgebiet}{ }

\renewcommand{\klausurtyp}{ }

\renewcommand{\klausurdatum}{ . 20}

\klausurvorspann {\fachbereich} {\klausurdatum} {\dozent} {\klausurgebiet} {\klausurtyp}

%Data untuk tabel nilai berikut

\renewcommand{\aeins}{  3  }

\renewcommand{\azwei}{  3   }

\renewcommand{\adrei}{  0  }

\renewcommand{\avier}{  0  }

\renewcommand{\afuenf}{  4  }

\renewcommand{\asechs}{  5  }

\renewcommand{\asieben}{  4  }

\renewcommand{\aacht}{  0  }

\renewcommand{\aneun}{  4  }

\renewcommand{\azehn}{  6  }

\renewcommand{\aelf}{  2  }

\renewcommand{\azwoelf}{  6  }

\renewcommand{\adreizehn}{  4   }

\renewcommand{\avierzehn}{  10  }

\renewcommand{\afuenfzehn}{  10  }

\renewcommand{\asechzehn}{  0  }

\renewcommand{\asiebzehn}{  61  }

\renewcommand{\aachtzehn}{    }

\renewcommand{\aneunzehn}{    }

\renewcommand{\azwanzig}{    }

\renewcommand{\aeinundzwanzig}{    }

\renewcommand{\azweiundzwanzig}{    }

\renewcommand{\adreiundzwanzig}{    }

\renewcommand{\avierundzwanzig}{    }

\renewcommand{\afuenfundzwanzig}{    }

\renewcommand{\asechsundzwanzig}{    }

\punktetabellesechzehn

\klausurnote

\newpage

\setcounter{section}{K}

\inputaufgabegibtloesung
{3}
{

Definisikan istilah-istilah berikut
\zusatzklammer {yang dicetak miring} {} {}.
\aufzaehlungsechs{Pemetaan
\stichwort {Weingarten} {}
pada $T_PY$ di suatu titik

\mavergleichskette
{\vergleichskette
{ P
}
{ \in }{ Y
}
{  }{
}
{    }{
}
{   }{
}
}
{}{}{}
dari suatu hiperpermukaan diferensiabel berorientasi

\mavergleichskette
{\vergleichskette
{ Y
}
{ \subseteq }{ \R^n
}
{  }{
}
{    }{
}
{   }{
}
}
{}{}{.}

}{Suatu
\stichwort {bundel vektor riil} {}
$V$ ber-rank $r$ di atas suatu
\definitionsverweis {ruang topologis}{}{}
$X$.

}{\stichwort {Integral lintasan} {} dari suatu bentuk diferensial-$1$
\mathl{\omega \in
{ \mathcal E }^{  1 } (  M   )}{} pada suatu manifold diferensiabel $M$ terhadap suatu kurva yang terdiferensialkan secara kontinu
\maabb {\gamma} {[a,b]} {M
} {.}

}{Suatu
\stichwort {bentuk volume positif} {}
$\omega$ pada suatu manifold diferensiabel $M$.

}{Suatu \stichwort {bentuk diferensial tertutup} {} $\omega$ pada suatu manifold diferensiabel $M$.

}{\stichwort {Percepatan tangensial} {}
dari suatu kurva yang dua kali terdiferensialkan secara kontinu
\maabb {\gamma} {I} {M
} {}
pada suatu
\definitionsverweis {manifold Riemann}{}{}
$M$.
}

}
{} {}

\inputaufgabegibtloesung
{3}
{

Rumuskan teorema-teorema berikut.
\aufzaehlungdrei{\stichwort {Teorema tentang kelengkungan normal} {.}}{Rumus untuk menghitung volume kanonik pada suatu manifold Riemann di dalam suatu peta.}{\stichwort {Teorema Stokes} {} untuk manifold dengan batas.}

}
{} {}

\inputaufgabe
{0}
{

}
{} {}

\inputaufgabe
{0}
{

}
{} {}

\inputaufgabegibtloesung
{4}
{

Buktikan teorema tentang hubungan antara kelengkungan dan vektor normal satuan dari suatu kurva bidang yang diparameterkan oleh panjang busur.

}
{} {}

\inputaufgabegibtloesung
{5 (1+2+2)}
{

Pada
\definitionsverweis {permukaan bola satuan}{}{}
$Y$ berdimensi dua, kita meninjau
\definitionsverweis {medan vektor tangensial}{}{}

\mavergleichskettedisp
{\vergleichskette
{ F \begin{pmatrix}  x  \\ y \\ z   \end{pmatrix}
}
{ =} { \begin{pmatrix}  y  \\ -x\\  0   \end{pmatrix}
}
{ } {
}
{ } {
}
{ } {
}
}
{}{}{.}
\aufzaehlungdrei{Tentukan
\definitionsverweis {matriks Jacobi}{}{}
dari $F$ pada $\R^3$.
}{Untuk titik

\mavergleichskette
{\vergleichskette
{ P
}
{ = }{ \begin{pmatrix}  0  \\ 1 \\  0   \end{pmatrix}
}
{ \in }{ Y
}
{    }{
}
{   }{
}
}
{}{}{}
tentukan suatu matriks untuk pemetaan
\maabbeledisp {} {T_PY} { T_PY
} { v } { \nabla_v F
} {,}
terhadap suatu
\definitionsverweis {basis}{}{}
yang sesuai.
}{Untuk titik

\mavergleichskette
{\vergleichskette
{ Q
}
{ = }{ \begin{pmatrix}  0  \\ 0\\  1   \end{pmatrix}
}
{ \in }{ Y
}
{    }{
}
{   }{
}
}
{}{}{}
tentukan suatu matriks untuk pemetaan
\maabbeledisp {} {T_QY} { T_QY
} { v } { \nabla_v F
} {,}
terhadap suatu
\definitionsverweis {basis}{}{}
yang sesuai.
}

}
{} {}

\inputaufgabegibtloesung
{4 (2+2)}
{

Kita meninjau elips

\mavergleichskettedisp
{\vergleichskette
{ E
}
{ =} { \left\{ (u,v)  \mid   2u^2 +5v^2 = 4         \right\}
}
{ } {
}
{ } {
}
{ } {
}
}
{}{}{.}
\aufzaehlungdrei{Definisikan suatu pemetaan surjektif yang terdiferensialkan secara kontinu
\maabb {} {\R} {E
} {.}
}{Deskripsikan suatu
\definitionsverweis {difeomorfisme}{}{}
antara lingkaran $S^1$ dan $E$.
}{
}

}
{} {}

\inputaufgabe
{0}
{

}
{} {}

\inputaufgabegibtloesung
{4}
{

Berikan suatu
\definitionsverweis {medan vektor}{}{} yang
\definitionsverweis {kontinu}{}{}
pada $S^2$ dan hanya mempunyai satu titik nol.

}
{} {}

\inputaufgabegibtloesung
{6}
{

Misalkan diberikan suatu
\definitionsverweis {data perekatan}{}{}

\mathbed {U_i} {}
 {i \in I} {}
 {} {} {} {,}
untuk
\definitionsverweis {ruang-ruang topologis}{}{.}
Buktikan bahwa terdapat suatu ruang topologis $X$ yang ditentukan secara unik, suatu penutup terbuka

\mavergleichskette
{\vergleichskette
{ X
}
{ = }{ \bigcup_{i \in I} V_i
}
{  }{
}
{    }{
}
{   }{
}
}
{}{}{}
dan
\definitionsverweis {homeomorfisme}{}{}
\maabb {\psi_i} {U_i } { V_i
} {}
sedemikian sehingga

\mavergleichskettedisp
{\vergleichskette
{ \psi_i  \left(  U_{ij}  \right)
}
{ =} { V_i \cap V_j
}
{ } {
}
{ } {
}
{ } {
}
}
{}{}{}
dan

\mavergleichskettedisp
{\vergleichskette
{ \psi_i \vert_{U_{ij} }
}
{ =} { \psi_j \vert_{U_{ji} }  \circ   \varphi_{ji}
}
{ } {
}
{ } {
}
{ } {
}
}
{}{}{}.

}
{} {}

\inputaufgabegibtloesung
{2}
{

Misalkan $M$ suatu
\definitionsverweis {manifold diferensiabel}{}{,}
\maabb {\gamma} { I } { M
} {}
suatu
\definitionsverweis {kurva diferensiabel}{}{}
di $M$, dan $\omega$ suatu
$1$-\definitionsverweis {bentuk}{}{}
terdiferensialkan pada $M$. Buktikan bahwa

\mavergleichskettedisp
{\vergleichskette
{ \gamma^* (d \omega)
}
{ =} { 0
}
{ } {
}
{ } {
}
{ } {
}
}
{}{}{}.

}
{} {}

\inputaufgabegibtloesung
{6 (3+3)}
{

Misalkan $M$ suatu
\definitionsverweis {manifold diferensiabel}{}{}
yang
\definitionsverweis {kompak}{}{}
dan mempunyai sedikitnya dua titik.
\aufzaehlungzwei {Buktikan bahwa suatu
$1$-\definitionsverweis {bentuk diferensial}{}{}
kontinu yang
\definitionsverweis {eksak}{}{}
pada $M$ mempunyai sedikitnya dua titik nol.
} {Buktikan bahwa hal ini tidak harus berlaku bagi suatu bentuk diferensial yang
\definitionsverweis {tertutup}{}{.}
}

}
{} {}

\inputaufgabegibtloesung
{4}
{

Berikan suatu contoh
\definitionsverweis {manifold dengan batas}{}{,}
yang
\definitionsverweis {terhubung}{}{}
dan batasnya terdiri atas
\zusatzklammer {gabungan saling lepas dari} {} {}
tak terhingga banyak salinan $\R^1$.

}
{} {}

\inputaufgabegibtloesung
{10}
{

Buktikan \stichwort {teorema tentang partisi kesatuan} {.}

}
{} {}

\inputaufgabegibtloesung
{10 (2+6+2)}
{

Kita meninjau
$2$-\definitionsverweis {bentuk diferensial}{}{}

\mavergleichskettedisp
{\vergleichskette
{ \omega
}
{ =} { x dy \wedge dz -y dx \wedge dz + z dx \wedge dy
}
{ } {
}
{ } {
}
{ } {
}
}
{}{}{}
pada bola satuan
\mathl{B  \left(  0 , 1  \right)}{.}
\aufzaehlungdreiabc{Buktikan bahwa $d \omega$ adalah tiga kali bentuk volume standar pada bola tersebut.
}{Buktikan bahwa
\mathl{\omega\vert_{S^2}}{} merupakan bentuk luas permukaan standar pada permukaan bola satuan.
}{Hitung luas permukaan bola dari
volume bola
dengan menggunakan
Teorema Stokes.
}

}
{} {}

\inputaufgabe
{0}
{

}
{} {}
